\section{\texorpdfstring{$\R$}{R}中区间的性质}

\begin{proposition}{$\R$中的闭区间都是闭集,开区间都是开集}{}
	\begin{enumerate}
		\item 对每个$a,b\in\R$,区间$\left[a,b\right]$是$\R$中的闭集,$\left(a,b\right)$是$\R$中的开集.
		\item 对每个$a\in\R$,区间$\left[a,+\infty\right),\left(-\infty,a\right]$是$\R$中的闭集,$\left(a,+\infty\right),\left(-\infty,a\right)$是$\R$中的开集.
	\end{enumerate}
\end{proposition}

\begin{proposition}{}{}
	$\left\{\left[-r,r\right]\subseteq\R\middle|r\in\Q_{>0}\right\}$是$\R$中原点的邻域基.
\end{proposition}

\begin{corollary}{$\R$的第一可数性}{}
	实直线是第一可数空间.
\end{corollary}

\begin{proposition}{(嵌入意义下)$\Q$在$\R$中稠密}{}
	设$\iota:\Q\to\R$是典范映射,则$\im\left(\iota\right)$在$\R$中稠密.
\end{proposition}

\begin{proposition}{$\R$中区间的序拓扑=$\R$诱导的子空间拓扑}{}
	设$I$是$\R$的区间,则$I$上的序拓扑与$\R$在$I$上诱导的子空间拓扑相同.
\end{proposition}

\begin{remark}{$\R$的第二可数性与可分性}{}
	如果$A$是$\R$的稠密子集,那么$\left\{\left(p,q\right)\subseteq\R\middle|p,q\in A\right\}$是$\R$的拓扑基,由此可知$\R$既是第二可数的又是可分的——取$A=\im\left(\iota\right)$便是,其中$\iota:\Q\to\R$是典范映射.
\end{remark}